\subsection{ARCH-136-GUARDED-PAIR: simultaneous individual primitives in a coefficient tail}
\label{sec:arch136guardedpair}

\paragraph{Scope.}
This is a positive lifting statement for a \emph{single specified}
coefficient multiple of an odd unit disk. It uses the horizontal
guard imposed in ARCH-136-DISK-SUPPORT. It does not select a family
of primitives compatible with all coefficient products or repeated
presentations. Those relations are not removed by this lemma.
The relevant sheets are cote 136, physical pp.~117, 119 and 121,
footers 116, 118 and 120,
\texttt{ARCHIVE\_136\_scan117.tex},
\texttt{ARCHIVE\_136\_scan119.tex}, and
\texttt{ARCHIVE\_136\_scan121.tex}.
Inputs are inherited AQ-136-COMPARISON and the relative row-extension
argument in AQ-136-MULTIPLICATIVE. The new cone and PD calculations
have tag none. No A, M, or outside preprint is used.

\paragraph{Statement.}
Let \(C\in\mathcal A_{\mathbf Z}\), let \(m\geq1\), let \(q\geq1\) be odd, and
suppose that the row \(C_*^{m,q}\) is horizontally acyclic,
including in degree zero. Fix an integer \(p\geq2\), and put
\[
                    N=p(q+1)-1.
\]
Choose \(W_{N+1}(C)\) for the inherited cohomology comparison
through form degree \(N+1\), uniformly in all simplicial sets.
For every \(s\geq0\) with
\[
                         pm+s\geq W_{N+1}(C),
 \tag{136.GP1}
\]
every simplicial set \(X\), and every
\(a\in C^{m,q}(X)\), there are homogeneous cochains
\[
 u\in C^{pm+s,N}(X),\qquad
 v\in C^{pm+s,N-1}(X)
\]
such that, with \(b=da\),
\[
 \begin{split}
 du&=T^{[s]}\gamma_p(b),\\
 dv&=p u-T^{[s]}a\gamma_{p-1}(b).
 \end{split}
 \tag{136.GP2}
\]
Thus a primitive for the divided-power cycle and a primitive for
its secondary relation exist simultaneously. For a finite family
of admissible targets one may take the maximum of their displayed
output-weight bounds. The choices themselves are not asserted
natural in all the data.

\paragraph{Proof.}
First extend \(a\) and zero from the two ends of \(X\times\Delta[1]\)
to a cochain
\[
        \widetilde a\in C^{m,q}(X\times\Delta[1]).
 \tag{136.GP3}
\]
This uses only the assumed acyclic row. Here is the precise lifting
argument, recalled from the relative part of
AQ-136-MULTIPLICATIVE. Under Dold--Kan, the normalized chains of
\(\mathbf Z[X\times\Delta[1]]/
  \mathbf Z[X\times\partial\Delta[1]]\)
are free, with the relative nondegenerate simplices as basis.
Extend the prescribed chain map to the normalized row one such
generator at a time in increasing degree. Its required boundary
has already been assigned and is a cycle. Acyclicity of the row
makes it a boundary and supplies a preimage. In degree zero an
arbitrary value can be chosen if no boundary is prescribed.
This constructs a simplicial module map, equivalently (136.GP3).
No linear contraction or projectivity of the terms of \(C\) is
assumed.

For a cochain complex \(A\), define the coefficient cone
\[
 P_p(A)^j=A^j\oplus A^{j+1},\qquad
 d_{P}(x,y)=(dx-py,-dy).
 \tag{136.GP4}
\]
Its cohomology has the natural exact sequence
\[
 0\longrightarrow H^j(A)/pH^j(A)
 \longrightarrow H^j(P_p(A))
 \longrightarrow H^{j+1}(A)[p]\longrightarrow0.
 \tag{136.GP5}
\]
This is the mapping-cone long exact sequence; it requires no
flatness of \(A\).

At weight \(w=pm+s\), consider the pair
\[
 F_s(a)=
 \bigl(T^{[s]}a\gamma_{p-1}(da),\
                         T^{[s]}\gamma_p(da)\bigr)
 \in P_p(C^{w,*}(X))^N .
 \tag{136.GP6}
\]
The special derivation and PD product identities show that it is
closed: the differential of its first entry is \(p\) times its
second entry, and its second entry is closed. The same construction
applied to \(\widetilde a\) gives a closed pair on
\(X\times\Delta[1]\), restricting to zero at one end and to
\(F_s(a)\) at the other.

Write \(\pi:X\times\Delta[1]\to X\) for projection. By (136.GP1)
and AQ-136-COMPARISON, its pullback induces isomorphisms on
cohomology of \(C^{w,*}\) in degrees \(N,N+1\): the corresponding
ordinary-cohomology projection is a homotopy equivalence.
The natural short exact sequences (136.GP5) and the short five
lemma therefore show that \(\pi^*\) also induces an isomorphism
on \(H^N(P_p(C^{w,*}))\). Both endpoint pullbacks are left
inverses of this isomorphism, and hence they are equal there.
The closed pair (136.GP6) must therefore have zero cohomology
class.

Choose a cone primitive \((r,t)\), so that
\[
 \bigl(T^{[s]}a\gamma_{p-1}(b),\
       T^{[s]}\gamma_p(b)\bigr)
              =(dr-pt,-dt).
 \tag{136.GP7}
\]
Setting \(u=-t\) and \(v=-r\) gives exactly (136.GP2), with the
asserted degrees and weights. The two bounds used, row acyclicity
at \((m,q)\) and comparison at \((pm+s,N+1)\), are independent
of \(X\). This proves the lemma. \(\square\)

\paragraph{Secondary-class consequence.}
With \(s=0\), once \(pm\geq W_{N+1}(C)\), the class
\(\Theta_p(da)\) of ARCH-136-PD-SECONDARY is defined and zero.
This is a statement for actual admissible theories, using their
horizontal row condition; it is compatible with the earlier
strict-PD example outside that category. For \(s>0\), (136.GP2)
concerns the coefficient multiple of the \emph{cycle and its
secondary relation}, rather than multiplying \(a\) before
applying a divided power. It is consequently different from
ARCH-136-PD-COEFFICIENT.

\paragraph{Remaining coherence gate.}
For each fixed \(p\), \(q\), guarded \(m\), and coefficient shift
satisfying (136.GP1), simultaneous target values for the two Moore
generators exist. This removes the secondary residue as an
obstruction to that individual tail lift. The proof does not
identify generators arising from different shifts or different
presentations, choose their values \(S\)-linearly, or prove
admissibility after adjoining all of them. In particular the
duplicate-pair phenomenon of ARCH-136-MOORE-DUPLICATE remains
relevant. The output-degree comparison bound is allowed to depend
on \(p\); no single prime-uniform completion is asserted.
